\documentclass[11pt, a4paper]{report}

% --- PAQUETES ESENCIALES Y GEOMETRÍA ---
\usepackage[a4paper, top=2.5cm, bottom=2.5cm, left=2.5cm, right=2.5cm]{geometry}
\usepackage{amsmath}
\usepackage{amssymb}
\usepackage{booktabs}
\usepackage{enumitem}
\usepackage{fancyhdr}
\usepackage{xcolor}

% --- PAQUETES DE TIKZ PARA EL DIAGRAMA ---
\usepackage{tikz}
% Corrección: Se cambia 'shadowbox' (que no existe) por la biblioteca 'shadows'
\usetikzlibrary{arrows.meta, positioning, shapes, fit, backgrounds, shadows}
% --- FIN PAQUETES DE TIKZ ---

% --- BLOQUE DE PREÁMBULO UNIVERSAL (FONT SPEC Y BABEL) ---
% Se eliminan inputenc, fontenc y lmodern, reemplazados por fontspec y babel
\usepackage{fontspec}
\usepackage[spanish, bidi=basic, provide=*]{babel}

\babelprovide[import, onchar=ids fonts]{spanish}
\babelprovide[import, onchar=ids fonts]{english}

% Establecer la fuente principal/latina (rm) a Noto Sans
\babelfont{rm}{Noto Sans}
% Como el español usa caracteres latinos, heredará Noto Sans.

% Corregir viñetas para idiomas no ingleses
% Se cambia \textbullet por - para compatibilidad
\setlist[itemize]{label=-}
% --- FIN DEL BLOQUE DE PREÁMBULO ---


% Color para la guía de uso
\definecolor{guidacolor}{rgb}{0.0, 0.2, 0.7} % Azul oscuro

% --- HYPERREF (CARGAR TARDE) ---
\usepackage{hyperref}
% Configuración de Hyperref
\hypersetup{
    colorlinks=true,
    linkcolor=black, % Links del ToC en negro por defecto
    urlcolor=cyan,
    citecolor=black
}

% Comandos personalizados
\newcommand{\destacado}[1]{\textbf{\textcolor{blue}{#1}}}
\newcommand{\advertencia}[1]{\par\vspace{0.5\baselineskip}\noindent\fcolorbox{red}{red!15}{\parbox{\dimexpr\linewidth-2\fboxsep-2\fboxrule\relax}{\textbf{ADVERTENCIA:} #1}}\par\vspace{0.5\baselineskip}}
\newcommand{\quotes}[1]{``#1''} % Comando para comillas tipográficas

% Configuración de encabezados y pies de página
\pagestyle{fancy}
\fancyhead{}
\fancyhead[R]{\nouppercase{\leftmark}}
\fancyhead[L]{Guía de Despliegue Data Guard}
\fancyfoot[C]{\thepage}
\renewcommand{\headrulewidth}{0.4pt}
\renewcommand{\footrulewidth}{0.4pt}

% Definición del estilo de listados de código para entornos Bash y SQL
\usepackage{listings}
\definecolor{codegray}{rgb}{0.5,0.5,0.5}
\definecolor{codeblue}{rgb}{0.0,0.0,0.6}
\lstset{
    basicstyle=\ttfamily\footnotesize, % Reducido de \small a \footnotesize
    breaklines=false,                 % Evita que las líneas se corten
    frame=single,
    showstringspaces=false,
    keywordstyle=\color{codeblue}\bfseries,
    commentstyle=\color{codegray},
    stringstyle=\color{red},
    numberstyle=\tiny\color{gray},
    language=bash
}

\lstdefinelanguage{SQL}{
  keywords={SELECT, FROM, WHERE, AND, OR, INSERT, UPDATE, DELETE, CREATE, DROP, ALTER, TABLE, VIEW, INDEX, AS, INTO, VALUES, COMMIT, ROLLBACK, BEGIN, END, SET, EXIT, 
ON, ALL, DATABASE, RECOVER, MANAGED, STANDBY, DISCONNECT, SESSION, LOGFILE, ARCHIVE, CURRENT, V, FROM, NOT, NULL, MAX, COUNT, ORDER, BY, GROUP, SHOW, PARAMETER, VDATABASE, VARCHIVEDLOG, VMANAGEDSTANDBY, VLOGFILE, VLOGHISTORY, VDATAGUARDSTATS, VARCHIVEGAP, VINSTANCE, DGMGRL, SYSTEM, REGISTER, LOG, UNTIL, SEQUENCE, THREAD, TO_CHAR, SYSDATE, USING, WITH, PDB, PLUGGABLE, ADD, DATAFILE, RESIZE, RUN, ALLOCATE, CHANNEL, FOR, MAINTENANCE, DEVICE, TYPE, DISK, DELETE, NOPROMPT, OBSOLETE, CROSSCHECK, LIST, BACKUP, SUMMARY, VALIDATE, RESTORE, RMAN, PLUS, DELETION, POLICY, RECOVERY, WINDOW, CONTROLFILE, AUTOBACKUP, FORMAT, RELEASE, REPORT, NEED, BLOCK, OFFLINE, ONLINE, SIZE, MEMBER},
  identifierstyle=\color{black},
  commentstyle=\color{codegray}\itshape,
  stringstyle=\color{red},
  basicstyle=\ttfamily\small,
  morecomment=[l]--
}


\title{\textbf{Manual de Instalación y Administración de Oracle Data Guard Automatizado}}
\author{
    Angelo Marin \\
    Joselyn Morales \\
    Kendall Salazar
}
\date{
    \textbf{Universidad Nacional de Costa Rica, Sede Región Brunca} \\
    \vspace{0.5em}
    Curso: Administración de Bases de Datos \\
    Profesor: Carlos Carranza \\
    \vspace{0.5em}
    Noviembre 2025
}

\begin{document}

\maketitle
\tableofcontents
\newpage

% *************************************************
% CAPÍTULO 1: INTRODUCCIÓN Y CONCEPTOS FUNDAMENTALES
% *************************************************

\chapter{Introducción y Conceptos Fundamentales}

\section{Despliegue Automatizado de Oracle Data Guard}

Este manual documenta la implementación completa de un sistema de alta disponibilidad utilizando Oracle Database 21c y Oracle Data Guard. El proceso está totalmente automatizado mediante scripts de Vagrant y Bash, permitiendo desplegar un entorno replicado (Principal y Standby) en minutos.

El sistema replica la base de datos principal en tiempo real a una base de datos \textit{standby} para garantizar la \textbf{cero pérdida de datos} y la \textbf{continuidad del servicio}. El enfoque de este manual está en cómo la arquitectura Data Guard automatizada resuelve la necesidad de \textbf{alta disponibilidad} para cualquier aplicación que lo requiera.

\subsection{Componentes Principales del Entorno}

\begin{enumerate}
    \item \textbf{Base de Datos Principal (\textit{Primary})}
    \begin{itemize}
        \item Oracle Database 21c Enterprise Edition.
\item Container Database (CDB): \texttt{cdb1}.
        \item Pluggable Database (PDB): \texttt{pdb1}.
        \item Esquema de aplicación: (Definido por scripts, \texttt{HC} por defecto).
        \item Servidor: \texttt{ol8-21-dg1.localdomain} (\texttt{192.168.56.101}).
\end{itemize}
    \item \textbf{Base de Datos \textit{Standby}}
    \begin{itemize}
        \item Réplica física exacta del \textit{primary}.
\item Database unique name: \texttt{cdb1\_stby}.
        \item Sincronización en tiempo real vía Oracle Data Guard.
        \item Servidor: \texttt{ol8-21-dg2.localdomain} (\texttt{192.168.56.102}).
\end{itemize}
    \item \textbf{Oracle Data Guard}
    \begin{itemize}
        \item Protección \texttt{MAXAVAILABILITY} (cero pérdida de datos).
\item Transporte \texttt{FASTSYNC} (sincrónico).
        \item Data Guard Broker (gestión centralizada).
        \item \textit{Switchover} y \textit{Failover} automáticos.
\end{itemize}
    \item \textbf{Infraestructura Automatizada}
    \begin{itemize}
        \item 2 VMs Oracle Linux 8 (VirtualBox).
\item Despliegue automatizado con Vagrant.
        \item Scripts PowerShell de orquestación.
        \item Backups automáticos diarios con RMAN.
\item Monitoreo continuo cada 10 minutos.
        \item Limpieza automática de \textit{archive logs}.
\end{itemize}
\end{enumerate}

\subsection{Objetivos Técnicos del Entorno}

El sistema automatizado está diseñado para cumplir con los objetivos técnicos estándar de un entorno de Alta Disponibilidad. Los scripts se encargan de:
\begin{itemize}
    \item \destacado{Desplegar dos servidores} (Primary y Standby) en VMs separadas.
    \item \destacado{Configurar la replicación} de Data Guard.
    \item \destacado{Establecer un \texttt{archive\_lag\_target}} (configurable en `install.env`, 300s por defecto) para forzar la conmutación de logs.
    \item \destacado{Automatizar la transferencia} y aplicación de logs, con scripts de monitoreo cada 10 minutos.
    \item \destacado{Implementar políticas de limpieza} de logs (configurable, 3 días por defecto).
    \item \destacado{Programar un respaldo diario automático} con RMAN y transferirlo al standby.
\end{itemize}

\subsection{Introducción a Oracle Data Guard}

Oracle Data Guard es la \textbf{solución de alta disponibilidad y recuperación ante desastres} de Oracle.
La filosofía de Data Guard se basa en copiar instantáneamente cada transacción del \textit{primary} a otro servidor. Si el primero falla, la copia toma el control, lo que se logra mediante la transferencia y aplicación de los \textbf{redo logs} y \textbf{archive logs}.

\subsubsection{Flujo de Replicación de Alto Nivel}

\begin{itemize}
    \item La \textbf{Primary Database} recibe todas las transacciones.
\item Oracle escribe los cambios en \textbf{redo logs}.
    \item Los \textit{redo logs} se convierten en \textbf{archive logs}.
\item Los \textit{archive logs} se \textbf{envían automáticamente al standby} vía red.
\item La \textbf{Standby database} aplica los \textit{archive logs} y se mantiene idéntica al \textit{primary}.
\end{itemize}

\section{Conceptos Fundamentales}

\subsection{Container Database (CDB) vs Pluggable Database (PDB)}

A partir de Oracle 12c, se utiliza la \textbf{arquitectura multitenant}:
\begin{itemize}
    \item \textbf{CDB (\textit{Container Database}):} Es el contenedor de nivel superior que aloja los metadatos del sistema y las PDBs. En este proyecto: \texttt{cdb1} y \texttt{cdb1\_stby}.
    \item \textbf{PDB (\textit{Pluggable Database}):} Es una base de datos autónoma dentro de un CDB. Contiene sus propios usuarios y esquemas (en este despliegue, los scripts cargan un esquema de ejemplo).
\end{itemize}
Esta arquitectura permite la consolidación de múltiples bases de datos.

\subsection{Primary Database vs Standby Database}

\begin{itemize}
    \item \textbf{Primary Database (Rol: Producción):} \textit{Read-Write}.
Recibe todas las transacciones de los usuarios. En este proyecto, \texttt{cdb1}.
    \item \textbf{Standby Database (Rol: Backup Activo):} \textit{Mounted} o \textit{Read-Only}.
Recibe replicación continua y aplica los cambios. En este proyecto, \texttt{cdb1\_stby}.
\end{itemize}

Ambas bases de datos tienen el mismo \texttt{DB\_NAME} (\texttt{cdb1}) porque son copias físicas, pero se distinguen por el \texttt{DB\_UNIQUE\_NAME} (\texttt{cdb1} y \texttt{cdb1\_stby}).
\subsection{Alta Disponibilidad (\textit{HA}) y Recuperación ante Desastres (\textit{DR})}

\begin{itemize}
    \item \textbf{Alta Disponibilidad (HA):} Minimiza el \textit{downtime}.
Las métricas clave son \textbf{RTO} (\textit{Recovery Time Objective}) y \textbf{RPO} (\textit{Recovery Point Objective}).
\item \textbf{Recuperación ante Desastres (DR):} Protege contra fallos catastficos.
\end{itemize}

Con Data Guard \texttt{MAXAVAILABILITY}, se garantiza un \textbf{RPO de cero segundos} (cero pérdida de datos) y un \textbf{RTO de minutos}.
\subsection{Replicación Sincrónica (\texttt{FASTSYNC})}

\begin{itemize}
    \item \textbf{Asincrónica (\texttt{ASYNC}):} \textit{Primary} no espera confirmación. Más rápido, pero posible pérdida de datos.
\item \textbf{Sincrónica (\texttt{SYNC} / \texttt{FASTSYNC}):} \textit{Primary} espera confirmación del \textit{standby}. Garantiza cero pérdida de datos.
\end{itemize}

El proyecto usa \texttt{FASTSYNC} para balancear la velocidad y la protección: el \textit{primary} espera que el \textit{standby} reciba los datos en memoria antes de continuar, lo cual es más rápido que esperar que los escriba en disco (\texttt{SYNC}) mientras mantiene el RPO en cero.
\subsection{Archive Logs}

Los \textbf{Archive Logs} son la copia histórica de los \textbf{Online Redo Logs} que contienen los registros de todos los cambios de la base de datos.
\begin{itemize}
    \item Son la base del mecanismo de replicación de Data Guard.
\item El parámetro \texttt{ARCHIVE\_LAG\_TARGET=300} (\textbf{5 minutos}) fuerza un \textit{log switch} periódico, asegurando que el \textit{standby} nunca esté más de 5 minutos desactualizado.
\end{itemize}

\subsection{Virtualización: VirtualBox y Vagrant}

\begin{itemize}
    \item \textbf{VirtualBox:} Es el \textbf{hypervisor} (motor de virtualización) que crea y ejecuta las Máquinas Virtuales (VMs) Oracle Linux 8.
    \item \textbf{Vagrant:} Es la herramienta de \textbf{orquestación} que lee el \texttt{Vagrantfile} y automatiza la creación, configuración y \textit{provisioning} de las VMs.
    \item \textbf{Provisioning (Aprovisionamiento):} En el contexto de Vagrant, este es el proceso de instalar software, configurar el sistema y ejecutar scripts automáticamente en la VM después de que esta ha sido creada. En este proyecto, los \textit{provisioners} son los scripts Bash que instalan Oracle y configuran Data Guard.
    \item \textbf{Synced Folders:} Permiten compartir directorios entre el \textit{host} (Windows) y las VMs (Linux), lo cual es **indispensable** para que los \textit{scripts} de \textit{provisioning} accedan al software de Oracle.
\end{itemize}

El Vagrantfile y Vagrant eliminan la necesidad de configurar manualmente las VMs, la red y el sistema operativo, siendo el comando clave \texttt{vagrant up}.

\subsection{Data Guard Broker (\texttt{DGMGRL})}
Este manual se basa exclusivamente en la herramienta \textbf{Data Guard Broker} (administrada mediante la línea de comandos \texttt{DGMGRL}) para todas las operaciones de administración.

\begin{itemize}
    \item \textbf{Gestión Centralizada:} El Broker es un framework que automatiza y centraliza la configuración y monitoreo de Data Guard.
    \item \textbf{Automatización de Tareas:} Simplifica tareas complejas como \textit{Switchover}, \textit{Failover} y \textit{Reinstatement} a un solo comando, en lugar de docenas de comandos SQL manuales.
    \item \textbf{Monitoreo (DMON):} Utiliza un proceso (DMON) para asegurar que la configuración esté sincronizada y reporta el estado de salud (\texttt{SUCCESS} o \texttt{ERROR}).
\end{itemize}
El uso del Broker es la práctica moderna recomendada por Oracle para administrar entornos Data Guard.

\section{Requisitos y Configuración Inicial}

Antes de iniciar el despliegue automatizado, es **necesario** verificar el cumplimiento de los siguientes requisitos.
\subsection{Requisitos de Hardware del Sistema Host (Su PC)}

Se recomienda un mínimo de \textbf{12 GB de RAM} para el sistema host, aunque \textbf{16 GB o más} son preferibles para un rendimiento fluido, ya que cada VM consumirá 4 GB.

Asegúrese de que la virtualización (Intel VT-x o AMD-V/SVM Mode) esté \textbf{Habilitada en la BIOS/UEFI} de su equipo.

\subsection{Requisitos de Software}

\begin{itemize}
    \item \textbf{Sistema Operativo:} Windows 10 o superior. El despliegue también puede ejecutarse en macOS o Linux, pero los scripts de orquestación de PowerShell (\texttt{deploy.ps1}) no funcionarán y deberán ejecutarse los comandos de \texttt{vagrant} manualmente.
    \item \textbf{VirtualBox:} Se recomienda la versión más reciente (ej. 7.0+ a Octubre 2025).
    \item \textbf{Vagrant:} Se recomienda la versión más reciente (ej. 2.4+ a Octubre 2025).
    \item \textbf{Software de Oracle:} El archivo \texttt{LINUX.X64\_213000\_db\_home.zip} ($\sim 2.9\text{ GB}$) debe ser descargado desde \texttt{oracle.com} y colocado en la carpeta \texttt{software/}. \textbf{No debe descomprimirse.}
\end{itemize}

\subsubsection{Instalación de VirtualBox y Vagrant}

Ambas herramientas deben instalarse siguiendo los pasos del instalador.

\advertencia{Se recomienda NO instalar el plugin \texttt{vagrant-vbguest}. Este plugin ha sido abandonado y puede causar conflictos con los scripts de \textit{provisioning} (específicamente, errores de PERL) al intentar recompilar las \textit{Guest Additions}.}

\subsection{Verificación y Configuración del Proyecto}

\subsubsection{Estructura de Carpetas Requerida}

Asegúrese de que existe el archivo del software de Oracle y los archivos de configuración en la estructura esperada:
\begin{itemize}
    \item \texttt{software/LINUX.X64\_213000\_db\_home.zip}
    \item \texttt{Vagrantfile}
    \item \texttt{config/vagrant.yml}
    \item \texttt{config/install.env}
    \item \texttt{powershell/deploy.ps1}
    \item Scripts en \texttt{2-Standby/scripts/}
\end{itemize}

\subsubsection{Configuración del Archivo \texttt{install.env}}

Este archivo contiene todas las variables de entorno usadas por los scripts Bash dentro de las VMs.
\begin{lstlisting}[language=bash, caption={Valores de configuración en install.env}]
# Passwords (cambiar en producción)
SYS_PASSWORD="S1sPassword1!"
PDB_PASSWORD="PdbPassword1!"
APP_SCHEMA_PASSWORD="password"
# IPs de Red (debe coincidir con vagrant.yml)
NODE1_PUBLIC_IP=192.168.56.101
NODE2_PUBLIC_IP=192.168.56.102

# Configuración de Data Guard
ARCHIVE_LAG_TARGET=300        # 5 minutos
PROTECTION_MODE=MAXAVAILABILITY
TRANSPORT_MODE=FASTSYNC

# Configuración de Backups
BACKUP_RETENTION_DAYS=7       # Retener backups 7 días
ARCHIVE_RETENTION_DAYS=3      # Retener archives 3 días
\end{lstlisting}

\advertencia{El archivo \texttt{install.env} contiene contraseñas en texto plano.
Asegúrese de que no se suba a repositorios públicos y de usar contraseñas fuertes en entornos de producción.}

\chapter{Arquitectura y Configuración Detallada}

\section{Arquitectura Completa del Sistema}

El sistema opera bajo una arquitectura de 3 capas (Presentación, Aplicación, Datos) consolidada en un entorno virtualizado.
\subsection{Componentes del Sistema}

\begin{enumerate}
    \item \textbf{Host Windows:} Ejecuta VirtualBox, Vagrant y orquestación con PowerShell.
\item \textbf{Primary Virtual Machine:} \texttt{ol8-21-dg1.localdomain} (\texttt{192.168.56.101}).
    \begin{itemize}
        \item Oracle Database 21c, CDB \texttt{cdb1}, PDB \texttt{pdb1}.
        \item Esquema de aplicación (por defecto \texttt{HC}).
\item Rol: \texttt{PRIMARY}, \textit{Read-Write}.
        \item Recursos: 4 GB RAM, 2 CPUs, 200 GB disco.
\end{itemize}
    \item \textbf{Standby Virtual Machine:} \texttt{ol8-21-dg2.localdomain} (\texttt{192.168.56.102}).
\begin{itemize}
        \item Oracle Database 21c, CDB \texttt{cdb1\_stby}, PDB \texttt{pdb1} (réplica).
\item Rol: \texttt{PHYSICAL STANDBY}, \textit{Read-Only}.
        \item Recursos: 4 GB RAM, 2 CPUs, 200 GB disco.
\end{itemize}
    \item \textbf{Data Guard:} Configuración \texttt{dg\_config} con \texttt{MAXAVAILABILITY} y \texttt{FASTSYNC}.
\end{enumerate}

\subsection{Diagrama de Arquitectura}
La arquitectura de red y almacenamiento se resume en el siguiente diagrama:

\begin{figure}[htbp]
  \centering
  \begin{tikzpicture}[
    node distance=1cm and 1.5cm,
    host/.style={rectangle, rounded corners, draw=black, fill=gray!10, minimum height=12cm, minimum width=15cm, label=above:\textbf{Host (Windows/Linux/macOS)}},
    % Corrección: Ancho reducido (min width 6cm -> 5.5cm)
    vm/.style={rectangle, rounded corners, draw=blue!80, fill=blue!5, minimum height=6.5cm, minimum width=5.5cm, thick, 
              drop shadow={opacity=0.7}},
    disk/.style={shape=cylinder, shape border rotate=90, draw=gray, fill=gray!20, aspect=0.3, minimum height=0.7cm, minimum width=2.5cm, font=\tiny},
    net/.style={rectangle, draw=green!60!black, dashed, fill=green!5, minimum width=14cm, font=\small},
    arrow/.style={-{Stealth[length=3mm]}, thick, draw=black!70},
    repl_arrow/.style={<->, Stealth-Stealth, thick, red!80, dashed}
  ]
    % Background Host Box
    \node (host) [host] {};

    % VM1 - Primary
    % Corrección: Aumentada separación horizontal (de -3.5 a -4.5)
    \node (vm1) [vm, at={(-4.5, 2.5)}] {};
    % Corrección: Texto anclado al NORTE (arriba)
    \node (vm1_text) [align=center, text width=5cm, anchor=north, at=(vm1.north), yshift=-0.2cm] {
        \textbf{VM1: Primary} \\
        \texttt{ol8-21-dg1} \\
        \texttt{192.168.56.101}
    };
    
    % Discos dentro de la VM
    \node (disk1_u01) [disk, at={(vm1.south)}, yshift=2.2cm] {\texttt{/u01} (Data/SW)};
    \node (disk1_u02) [disk, below=0.2cm of disk1_u01] {\texttt{/u02} (FRA/Backup)};
    
    % VM2 - Standby
    % Corrección: Aumentada separación horizontal (de 3.5 a 4.5)
    \node (vm2) [vm, at={(4.5, 2.5)}] {};
    % Corrección: Texto anclado al NORTE (arriba)
    \node (vm2_text) [align=center, text width=5cm, anchor=north, at=(vm2.north), yshift=-0.2cm] {
        \textbf{VM2: Standby} \\
        \texttt{ol8-21-dg2} \\
        \texttt{192.168.56.102}
    };

    % Discos dentro de la VM
    \node (disk2_u01) [disk, at={(vm2.south)}, yshift=2.2cm] {\texttt{/u01} (Data/SW)};
    \node (disk2_u02) [disk, below=0.2cm of disk2_u01] {\texttt{/u02} (FRA/Backup)};
    
    % Redes
    \node (net_hostonly) [net, below=1cm of host.south, anchor=center] {\textbf{Red Host-Only (192.168.56.0/24)}};
    \node (net_nat) [net, above=1cm of host.north, anchor=center] {\textbf{Red NAT (Acceso a Internet)}};

    % Conexiones
    \draw [arrow] (vm1.south) -- (vm1.south |- net_hostonly.north);
    \draw [arrow] (vm2.south) -- (vm2.south |- net_hostonly.north);
    
    % Corrección: Texto "Replicación" subido (above=3pt)
    \draw [repl_arrow] (vm1.east) -- (vm2.west) node [midway, above=3pt, fill=white, rounded corners, font=\small] {Replicación Data Guard};
    
    \draw [arrow] (vm1.north) -- (vm1.north |- net_nat.south);
    \draw [arrow] (vm2.north) -- (vm2.north |- net_nat.south);
    
  \end{tikzpicture}
  \caption{Diagrama de la arquitectura de despliegue (TikZ).}
  \label{fig:arquitectura_tikz}
\end{figure}

\subsection{Flujo de Replicación de Datos Data Guard}

\begin{enumerate}
    \item Se produce una transacción en \textbf{Primary}.
\item El proceso \texttt{LGWR} escribe el cambio en \textit{Redo Logs}.
    \item El proceso \texttt{ARCn} crea un \textit{Archive Log}.
\item El proceso \texttt{LNS} (\textit{Log Network Server}) envía el \textit{redo data} al \textit{Standby} usando \texttt{FASTSYNC} sobre \texttt{192.168.56.0/24}.
\item El proceso \texttt{RFS} (\textit{Remote File Server}) en el \textit{Standby} recibe los datos.
\item El proceso \texttt{MRP} (\textit{Managed Recovery Process}) en el \textit{Standby} aplica los cambios a los \textit{datafiles} del \textit{standby}.
\end{enumerate}

\section{Orquestación de Máquinas Virtuales (\texttt{Vagrantfile})}

El \texttt{Vagrantfile} define la infraestructura como código.
\subsection{Estructura y Definición de VMs}

El archivo define dos bloques, uno para \texttt{primary} y uno para \texttt{standby}, utilizando la configuración cargada desde \texttt{config/vagrant.yml}.
\begin{lstlisting}[language=ruby, caption={Estructura del Vagrantfile}]
Vagrant.configure("2") do |config|
  # Carga de config externa
  params = YAML.load_file 'config/vagrant.yml'
  config.vm.box = params['shared']['box']
  # ... Synced folders ...

  # Definición de Primary VM
  config.vm.define "primary", primary: true do |primary|
primary.vm.hostname = "ol8-21-dg1.localdomain"
    primary.vm.network "private_network", ip: params['node1']['public_ip']
    primary.vm.network "forwarded_port", guest: 1521, host: 1521
    # ... Configuracion de recursos y discos (VirtualBox Provider) ...
    primary.vm.provision "shell", path: "2-Standby/scripts/primary/setup.sh"
  end

  # Definición de Standby VM
  config.vm.define "standby", autostart: false do |standby|
standby.vm.hostname = "ol8-21-dg2.localdomain"
    standby.vm.network "private_network", ip: params['node2']['public_ip']
    standby.vm.network "forwarded_port", guest: 1521, host: 1522 % Puerto diferente
    # ...
    standby.vm.provision "shell", path: "2-Standby/scripts/standby/setup.sh"
  end
\end{lstlisting}

\subsection{Carpetas Compartidas (\textit{Synced Folders})}

Las carpetas compartidas son un componente principal del \textit{provisioning}:
\begin{itemize}
    \item \texttt{.} (\textit{host root}) $\to$ \texttt{/vagrant} (\textit{VM})
    \item \texttt{software/} $\to$ \texttt{/vagrant\_software} (\textbf{Software de Oracle})
    \item \texttt{config/} $\to$ \texttt{/vagrant\_config} (\textbf{Archivos \texttt{install.env}} y \texttt{vagrant.yml})
\end{itemize}

\subsection{Configuración de Vagrant (\texttt{vagrant.yml} e \texttt{install.env})}

\subsubsection{\texttt{config/vagrant.yml} (YAML)}

Controla los recursos de infraestructura y es leído por el \texttt{Vagrantfile}.
\begin{lstlisting}[language=, caption={Extracto de vagrant.yml}]
node1:
  vm_name: ol8_21_dg1_primary
  mem_size: 4096       # 4 GB de RAM
  cpus: 2              # 2 Cores
  public_ip: 192.168.56.101
  disk_size: 102400    # 100 GB (para u01 y u02)
\end{lstlisting}

\subsubsection{\texttt{config/install.env} (Bash)}

Define las variables de entorno de la aplicación Oracle y Data Guard, leído por los \textit{scripts} Bash:

\begin{lstlisting}[language=bash, caption={Valores de configuración en install.env}]
# Passwords (cambiar en producción)
SYS_PASSWORD="S1sPassword1!"
PDB_PASSWORD="PdbPassword1!"
APP_SCHEMA_PASSWORD="password"

# IPs de Red (debe coincidir con vagrant.yml)
NODE1_PUBLIC_IP=192.168.56.101
NODE2_PUBLIC_IP=192.168.56.102

# Configuración de Data Guard
ARCHIVE_LAG_TARGET=300        # 5 minutos
PROTECTION_MODE=MAXAVAILABILITY
TRANSPORT_MODE=FASTSYNC

# Configuración de Backups
BACKUP_RETENTION_DAYS=7       # Retener backups 7 días
ARCHIVE_RETENTION_DAYS=3      # Retener archives 3 días
\end{lstlisting}

Para modificar la configuración del ambiente (RAM, IPs, \textit{passwords}), es necesario editar estos archivos y luego ejecutar \texttt{.\textbackslash powershell\textbackslash destroy.ps1} seguido de \texttt{.\textbackslash powershell\textbackslash deploy.ps1} para aplicar los cambios.

\section{Orquestación y Compatibilidad}

\subsection{Script de Despliegue Principal (\texttt{deploy.ps1})}

El \texttt{deploy.ps1} (PowerShell) es el orquestador principal que ejecuta todo el proceso en 6 fases:

\begin{enumerate}
    \item 
\textbf{Verificación de Prerequisitos:} Comprueba software y recursos.
    \item \textbf{Corrección de \textit{Line Endings}:} Ejecuta \texttt{fix-line-endings.ps1}.
\item \textbf{Limpieza de VMs Antiguas:} Ejecuta \texttt{vagrant destroy -f} y \texttt{VBoxManage unregistervm --delete}.
\item \textbf{Despliegue de \textit{Primary}:} Ejecuta \texttt{vagrant up primary} (30-45 minutos).
    \item \textbf{Despliegue de \textit{Standby}:} Ejecuta \texttt{vagrant up standby} (20-30 minutos).
\item \textbf{Verificación Post-Despliegue y Logs:} Verifica SSH, \textit{cron jobs} y recolecta \textit{alert logs} de Oracle.
\end{enumerate}

\subsection{Script de Corrección de \textit{Line Endings} (\texttt{fix-line-endings.ps1})}

Este script es \textbf{importante} para la compatibilidad entre Windows y Linux.
El problema $\texttt{\$'\textbackslash r': command not found}$ es el síntoma de \textit{line endings} incorrectos.
El script convierte los archivos \texttt{.sh} y \texttt{.env} de formato Windows \textbf{CRLF} (\texttt{\textbackslash r\textbackslash n}) a formato Unix \textbf{LF} (\texttt{\textbackslash n}) para que Bash pueda interpretarlos correctamente.

\begin{lstlisting}[language=bash, caption={Ejecución de fix-line-endings.ps1}]
# El script se ejecuta automáticamente en la Fase 2
& powershell -ExecutionPolicy Bypass -File .\powershell\fix-line-endings.ps1
\end{lstlisting}
Esta acción previene el fallo de \textbf{todos} los \textit{provisioners} Bash.

% Cambia el color del link en el ToC
\addtocontents{toc}{\protect\hypersetup{linkcolor=guidacolor}}
\chapter{\textcolor{guidacolor}{Guía de Uso y Operaciones}}
% Restaura el color original del link en el ToC
\addtocontents{toc}{\protect\hypersetup{linkcolor=black}}

\section{Despliegue y Uso del Sistema}

Antes de ejecutar cualquier comando, asegúrese de haber cumplido todos los requisitos de software listados en la Sección 1.8, incluyendo la descarga del archivo \texttt{LINUX.X64\_213000\_db\_home.zip} en la carpeta \texttt{software/}.

\subsection{Despliegue Completo Automatizado}

La manera más eficiente de desplegar el sistema es a través del \textit{script} orquestador:

\begin{lstlisting}[language=bash, caption={Comando Principal de Despliegue}]
.\powershell\deploy.ps1
\end{lstlisting}

Este comando se encargará de instalar todo, desde el sistema operativo en las VMs hasta la configuración final de Data Guard Broker, incluyendo la creación del \textit{standby} con \quotes{...RMAN DUPLICATE FOR STANDBY FROM ACTIVE DATABASE.}

\subsection{Operaciones Individuales con Vagrant}

\begin{itemize}
    \item \textbf{Conexión SSH:} \texttt{vagrant ssh primary} o \texttt{vagrant ssh standby}.
\item \textbf{Apagar:} \texttt{vagrant halt primary} o \texttt{vagrant halt standby}.
    \item \textbf{Reanudar:} \texttt{vagrant up primary} o \texttt{vagrant up standby} (si ya existen).
\item \textbf{Re-ejecutar \textit{Provisioners}:} \texttt{vagrant provision primary} (aplicar cambios en \textit{scripts} Bash).
\item \textbf{Configurar Auto-Start de Oracle:} \texttt{vagrant provision primary --provision-with setup\_oracle\_autostart}.
\end{itemize}

\subsection{Destrucción del Ambiente}

\begin{lstlisting}[language=bash, caption={Destrucción Completa (Recomendado)}]
.\powershell\destroy.ps1
\end{lstlisting}
Este comando garantiza que se eliminen completamente las VMs, los discos virtuales y los registros de VirtualBox, liberando espacio en disco ($\sim 200\text{ GB}$).
\section{Conexiones a las Bases de Datos}

\subsection{Conexiones desde el Host Windows}

\begin{table}[htbp]
    \centering
    \caption{Credenciales y Puertos de Conexión}
    \label{tab:conexiones_host}
    \begin{tabular}{lll}
        \toprule
        \textbf{Rol / Usuario} & \textbf{Conexión (Puerto Host)} & \textbf{Contraseña} \\
        \midrule
        \textbf{Primary (SYSDBA)} & \texttt{sys/S1sPassword1!@localhost:1521/cdb1} & \texttt{S1sPassword1!} \\
        \textbf{Primary (Schema)} & \texttt{HC/password@localhost:1521/pdb1} & \texttt{password} \\
        \textbf{Standby (SYSDBA)} 
& \texttt{sys/S1sPassword1!@localhost:1522/cdb1\_stby} & \texttt{S1sPassword1!} \\
        \bottomrule
    \end{tabular}
\end{table}

Recuerde que el puerto para el \textit{Standby} desde el \textit{Host} es \texttt{1522}, debido al reenvío de puertos configurado en el \texttt{Vagrantfile}.

\subsubsection{Conexión con SQL Developer (o GUI similar)}
Para conectarse a la Pluggable Database (PDB) donde reside el esquema de aplicación (por defecto \texttt{HC}) usando una herramienta gráfica, es importante usar la configuración de red correcta.

\begin{itemize}
    \item \textbf{Username:} \texttt{HC}
    \item \textbf{Password:} \texttt{password} (o el valor de \texttt{APP\_SCHEMA\_PASSWORD})
    \item \textbf{Hostname:} \texttt{localhost}
    \item \textbf{Port:} \texttt{1521} (para el Primary)
    \item \textbf{Connection Type:} Seleccione \textbf{Service name} (No usar SID).
    \item \textbf{Service name:} \texttt{pdb1.world}
\end{itemize}
El \textit{service name} \texttt{pdb1.world} es el nombre global de la base de datos PDB, que es requerido para la conexión \textit{multitenant}.

\subsection{Scripts de Verificación del Sistema}

El script más importante para el monitoreo post-despliegue es:

\begin{lstlisting}[language=bash, caption={Verificación Completa de Data Guard}]
vagrant ssh primary -c "/home/oracle/automation/verify_dataguard.sh"
\end{lstlisting}
Este script produce un reporte consolidado con el rol de la base de datos, el estado de los procesos (\texttt{MRP} / \texttt{LNS}), el \textit{apply lag} (debe ser $\mathbf{0\text{ segundos}}$), y la verificación de la replicación del esquema de aplicación.
\section{Scripts de Mantenimiento y Automatización}

Todos los scripts de automatización se encuentran en \texttt{/home/oracle/automation/} dentro de las VMs y se ejecutan bajo el usuario \texttt{oracle}.
\subsection{Scripts Cron Automatizados}

\begin{itemize}
    \item \textbf{\texttt{check\_log\_transport.sh} (Cada 10 min en \textit{Primary}):} Verifica el \textit{lag} de aplicación entre el \textit{primary} y el \textit{standby}.
\item \textbf{\texttt{rman\_daily\_backup.sh} (Diario 1 AM en \textit{Primary}):} Ejecuta \texttt{RMAN BACKUP DATABASE PLUS ARCHIVELOG} con compresión, elimina obsoletos (\textit{retention policy} de \textbf{7 días}) y transfiere la copia al \textit{standby} vía \texttt{rsync} sobre SSH.
\item \textbf{\texttt{delete\_archives\_primary.sh} (Diario 2 AM en \textit{Primary}):} Elimina \textit{archive logs} que tienen más de \textbf{3 días} y han sido respaldados o aplicados en el \textit{standby}.
\item \textbf{\texttt{delete\_archives\_standby.sh} (Diario 11 PM en \textit{Standby}):} Elimina \textit{archive logs} que tienen más de \textbf{3 días} y han sido aplicados.
\end{itemize}

\subsection{Scripts de Operación Manual}

\begin{itemize}
    \item \textbf{\texttt{manual\_backup.sh}:} Ejecución de \textit{backup} RMAN a demanda antes de operaciones críticas.
    \item \textbf{\texttt{manual\_sync.sh}:} Fuerza un \textit{log switch} inmediato para verificar la sincronización o reducir el \textit{lag} tras una actividad intensa.
\end{itemize}

\section{Troubleshooting Rápido}

\begin{itemize}
    \item \textbf{Problema:} Deploy falla con \texttt{\$'\textbackslash r': command not found}.
\item \textbf{Solución:} Ejecutar manualmente \texttt{.\textbackslash powershell\textbackslash fix-line-endings.ps1} y re-deployar.
    \item \textbf{Problema:} \texttt{ORA-12541: TNS:no listener} o \texttt{ORA-01034: ORACLE not available}.
\item \textbf{Solución:} Iniciar manually el \textit{listener} (\texttt{lsnrctl start}) y la base de datos (\texttt{STARTUP}) a través de SSH. Si el problema es recurrente, configure el \textit{auto-start} con el \textit{provisioner} \texttt{setup\_oracle\_autostart}.
    \item \textbf{Problema:} \textit{Apply lag} en \textit{Standby} aumenta.
\item \textbf{Solución:} El proceso \texttt{MRP0} del \textit{Standby} está lento o detenido. Verifique su estado con \texttt{v\$managed\_standby} y, si es necesario, reinicie la aplicación de logs (\texttt{ALTER DATABASE RECOVER MANAGED STANDBY DATABASE DISCONNECT FROM SESSION}).
\end{itemize}

Ante cualquier fallo irrecuperable, la solución más rápida es \textbf{destruir y re-deployar} con el script \texttt{deploy.ps1}, asegurando que la configuración en \texttt{install.env} y \texttt{vagrant.yml} sea correcta.

\newpage

% *************************************************
% CAPÍTULO 8: OPERACIONES Y MANTENIMIENTO DIARIO (90-94)
% *************************************************

\chapter{Operaciones y Mantenimiento Diario}

Esta sección documenta los procedimientos esenciales que el DBA debe ejecutar para mantener la salud del sistema, incluyendo el monitoreo, \textit{switchover}, \textit{failover} y \textit{troubleshooting}.
\section{Monitoreo Diario de Data Guard}

\subsection{Procedimientos Automatizados}
\begin{itemize}
    \item \textbf{Monitoreo de Lag:} \texttt{check\_log\_transport.sh} (cada 10 min) verifica que el \texttt{Apply Lag} sea $\mathbf{< 5 \text{ segundos}}$.
\item \textbf{Backups y Limpieza:} \texttt{rman\_daily\_backup.sh} (1 AM) y \texttt{delete\_archives\_primary.sh} (2 AM) se ejecutan automáticamente.
\end{itemize}

\subsection{Verificación Manual (\texttt{verify\_dataguard.sh})}
El DBA debe ejecutar este reporte para confirmar el estado:
\begin{lstlisting}[language=bash, caption={Uso del Script de Verificación}]
vagrant ssh primary -c "/home/oracle/automation/verify_dataguard.sh"
\end{lstlisting}
\destacado{Puntos Clave a Revisar:} \texttt{Configuration Status} debe ser \texttt{SUCCESS}, \texttt{Apply Lag} debe ser mínimo, y \texttt{Archive Gap} debe ser \texttt{no rows selected}.
\subsection{Monitoreo de Recursos}
\begin{itemize}
    \item \textbf{Espacio en Disco (\texttt{df -h}):} \texttt{/u01} y \texttt{/u02} deben estar bajo el 85\% de uso.
\item \textbf{Logs:} Revisar los \textit{alert logs} del \textit{primary} y \textit{standby} en \texttt{/u01/app/oracle/diag/} por errores \texttt{ORA-}.
\end{itemize}

\section{Respaldos Manuales y Verificación}

\subsection{Ejecución On-Demand (\texttt{manual\_backup.sh})}
\begin{enumerate}
    \item Ejecutar \texttt{/home/oracle/automation/manual\_backup.sh} en el \textit{Primary}.
\item Confirmar \textit{backup} y transferencia al \textit{Standby}.
\end{enumerate}

\subsection{Verificación de Backups}
\begin{itemize}
    \item \textbf{RMAN:} Ejecutar \texttt{LIST BACKUP SUMMARY} para confirmar la existencia.
\item \textbf{Validación:} Ejecutar \texttt{RESTORE DATABASE VALIDATE} para simular una restauración sin afectar la base de datos.
\end{itemize}

\section{Switchover Planificado}

\subsection{Procedimiento de Cambio de Roles}
\begin{enumerate}
    \item \textbf{Pre-Check:} Ejecutar \texttt{verify\_dataguard.sh} para confirmar \texttt{Apply Lag} mínimo y \texttt{Switchover Status = TO STANDBY}.
\item \textbf{Ejecutar Switchover:}
        \begin{lstlisting}[language=bash, caption={Comando de Switchover con DGMGRL}]
dgmgrl /
SWITCHOVER TO cdb1_stby;
\end{lstlisting}
    \item \textbf{Verificación Post-Switchover:} \texttt{cdb1\_stby} se convierte en \texttt{PRIMARY} (READ WRITE) y \texttt{cdb1} se convierte en \texttt{PHYSICAL STANDBY} (APPLY-ON).
\item \destacado{Resultado:} Downtime de $\sim 2$-$5$ minutos, cero pérdida de datos.
\end{enumerate}

\section{Failover de Emergencia}

\subsection{Procedimiento de Activación}
\begin{enumerate}
    \item \textbf{Evaluación:} Confirmar que el \textit{Primary} está inoperable (\textit{no responde SSH/TNS}).
\item \textbf{Ejecutar Failover:}
        \begin{lstlisting}[language=bash, caption={Comando de Failover de Emergencia}]
dgmgrl sys/password@cdb1_stby
FAILOVER TO cdb1_stby;
\end{lstlisting}
    \item \textbf{Activación:} El \textit{Standby} se activa como \texttt{PRIMARY} (\textit{READ WRITE}).
\item \textbf{Pérdida de Datos:} \texttt{MAXAVAILABILITY} asegura que la pérdida sea mínima (o cero) si el \textit{Standby} estaba sincronizado.
\end{enumerate}

\subsection{Reinstatement del Primary Fallido}
\begin{enumerate}
    \item \textbf{Recuperar Host:} Reiniciar la VM del \textit{Primary} caído (\texttt{vagrant reload primary}).
\item \textbf{Reinstaurar:} Utilizar \textit{Broker} para convertir la base de datos fallida en el nuevo \textit{Standby}:
        \begin{lstlisting}[language=bash, caption={Comando de Reinstatement (Rápido con Flashback)}]
dgmgrl /
REINSTATE DATABASE cdb1;
\end{lstlisting}
\advertencia{El comando \texttt{REINSTATE} depende de que \textbf{Flashback Database} esté activo y que el tiempo de retención de los \textit{flashback logs} no haya sido excedido. Si falla, se requerirá una recreación manual del standby.}
\end{enumerate}
\destacado{Resultado:} El \textit{old primary} se convierte en \textit{standby} en $\sim 5$ minutos gracias a \textit{Flashback Database}.
\section{Troubleshooting General}

\subsection{Matriz de Problemas Comunes}

\begingroup
\setlength{\tabcolsep}{6pt} % Espacio entre columnas
\par\noindent
\begin{tabular*}{\textwidth}{@{\extracolsep{\fill}}lll}
    \toprule
    \textbf{Síntoma} & \textbf{Causa Probable} & \textbf{Solución Inmediata} \\
    \midrule
    \texttt{Apply lag} alto & \texttt{MRP0} detenido / I/O lento & Iniciar \texttt{MRP0} y revisar \textit{Alert Log} \\
    \addlinespace
    \texttt{Archive gap} detectado & Logs perdidos o \textit{Primary} caído & Copiar \textit{archives} manualmente y registrar \\
    \addlinespace
    \texttt{ORA-12541} & \textit{Listener} caído & \texttt{lsnrctl start} \\
    \addlinespace
    \texttt{Disk full} (\texttt{FRA}) & Archives acumulados & \texttt{DELETE NOPMPT OBSOLETE} \\
    \addlinespace
    \texttt{SSH connection} falla & \textit{Host Keys} no aceptadas & Ejecutar \texttt{setup\_ssh\_keys.sh} \\
    \bottomrule
\end{tabular*}
\par\vspace{\baselineskip} % Añade espacio vertical después de la tabla
\endgroup

Para la mayoría de los errores, la primera acción es revisar el \textbf{\textit{Alert Log}} del nodo (\texttt{tail -f /u01/.../alert\_cdb1.log}) para identificar el código \texttt{ORA-XXXXX} específico.

\newpage
% *************************************************
% CAPÍTULO 9: CONCLUSIONES Y RECOMENDACIONES (99)
% *************************************************

\chapter{Conclusiones y Recomendaciones Finales}

\section{Resumen Ejecutivo y Logros}

Este proyecto implementa una solución de \textbf{Alta Disponibilidad (HA)} y \textbf{Disaster Recovery (DR)} completa y automatizada.
\subsection{Logros del Proyecto}
\begin{itemize}
    \item \textbf{Zero Data Loss (RPO=0):} Asegurado por \texttt{MAXAVAILABILITY} y transporte \texttt{FASTSYNC}.
\item \textbf{Deployment Automatizado:} Configuración completa en $\sim 60$-$70$ minutos con \texttt{.\textbackslash powershell\textbackslash deploy.ps1}.
\item \textbf{Recuperación Rápida:} \textit{Failover} en minutos, con \textit{reinstatement} del \textit{primary} en $\sim 5$ minutos (\textit{Flashback Database}).
\item \textbf{Protección por Capas:} Replicación, \textit{backups} RMAN diarios y copia \textit{off-site} (\texttt{rsync} a \textit{Standby}).
\item \textbf{Monitoreo Proactivo:} \textit{Cron jobs} cada 10 minutos para verificar \textit{apply lag} y \textit{archive gaps}.
\item \textbf{Documentación:} Guía exhaustiva lista para auditoría y transferencia de conocimiento.
\end{itemize}

\section{Recomendaciones para Producción}

Si esta arquitectura se traslada a un entorno de producción:

\subsection{Infraestructura y Red}
\begin{itemize}
    \item \textbf{Hardware Dedicado:} Utilizar servidores físicos o instancias IaaS (\texttt{OCI}, \texttt{AWS}) con \textbf{mínimo 32 GB de RAM} y \textbf{dispositivos SSD empresariales} para el almacenamiento.
\item \textbf{Red Geográfica:} Colocar los nodos \textit{Primary} y \textit{Standby} en \textbf{centros de datos separados} para protección real ante desastres regionales.
\item \textbf{Conectividad:} Garantizar \textbf{1 Gbps de ancho de banda dedicado} y baja latencia para la replicación \texttt{FASTSYNC}.
\end{itemize}

\subsection{Operación y Seguridad}
\begin{itemize}
    \item \textbf{Monitoreo Enterprise:} Reemplazar scripts Bash por \textbf{Oracle Enterprise Manager (OEM)} o herramientas como \textbf{Prometheus/Grafana} para alertas automatizadas (email/SMS).
\item \textbf{TDE:} Implementar \textbf{Transparent Data Encryption (TDE)} para cifrar los \textit{datafiles} de la base de datos en reposo.
\item \textbf{Pruebas de DR:} Realizar ejercicios de \textit{Switchover} y \textit{Failover} al menos \textbf{trimestralmente} para validar procedimientos y medir el \texttt{RTO} real.
\item \textbf{Licenciamiento:} Verificar la necesidad de licencias adicionales (\textit{e.g.}, \textit{Active Data Guard} para \textit{queries} en el \textit{standby}).
\end{itemize}

\section{Posibles Mejoras Futuras}

\begin{itemize}
    \item \textbf{Failover Automático (\texttt{FSFO}):} Habilitar \texttt{Fast-Start Failover} con un proceso \textit{Observer} para eliminar la dependencia de la intervención del DBA en caso de fallo.
\item \textbf{Clonación Rápida:} Utilizar \texttt{SNAPSHOT STANDBY} para crear copias \textit{Read-Write} temporales para QA y desarrollo a partir del \textit{standby}.
\item \textbf{Automatización DevOps:} Integración con \texttt{Ansible} o \texttt{GitLab CI} para la gestión de configuración y el despliegue de parches.
\end{itemize}

\section{Conclusión Final}

La implementación de este proyecto demuestra la notable madurez y eficiencia de las herramientas modernas de Oracle para la alta disponibilidad. La combinación de Oracle Data Guard 21c y el Data Guard Broker ha simplificado drásticamente lo que históricamente era un proceso complejo y propenso a errores.

El comando \texttt{PREPARE DATABASE FOR DATA GUARD} del Broker, por ejemplo, encapsula una multitud de pasos manuales (como la activación de \texttt{FORCE LOGGING}, modo \texttt{ARCHIVELOG} y la creación de \textit{Standby Redo Logs}) en una sola operación robusta.

Asimismo, la capacidad de RMAN para clonar una base de datos completa a través de la red con \texttt{DUPLICATE TARGET DATABASE FOR STANDBY FROM ACTIVE DATABASE} elimina la necesidad de realizar respaldos y restauraciones manuales, construyendo la réplica \textit{standby} de forma automatizada y fiable.

Finalmente, la gestión del ciclo de vida de la replicación (operaciones de \textit{switchover} y \textit{failover}) se reduce a comandos únicos y controlados. La capacidad de reinstaurar un \textit{primary} fallido como un nuevo \textit{standby} usando \texttt{REINSTATE DATABASE}, aprovechando \textit{Flashback Database}, subraya la resiliencia y facilidad de recuperación del sistema.

Si bien la automatización de la infraestructura con herramientas como Vagrant y VirtualBox fue un facilitador clave para este entorno de demostración, el verdadero logro reside en la potencia y simplicidad que las propias herramientas de Oracle Data Guard proporcionan para garantizar la integridad de los datos y la continuidad del servicio.


\newpage
\appendix

% *************************************************
% APÉNDICE A: FLUJO DE DESPLIEGUE DEL NODO PRIMARY
% *************************************************

\chapter{Análisis Detallado: Despliegue del Nodo Principal}

Esta sección documenta detalladamente el proceso de \textit{provisioning} del nodo \textit{primary} (\texttt{ol8-21-dg1}).
El flujo se inicia con \texttt{vagrant up primary} que llama al \textit{wrapper} \texttt{setup.sh} y se orquesta completamente en \texttt{root\_setup.sh}.
\section{Punto de Entrada (\texttt{setup.sh})}

El script \texttt{setup.sh} es el primer script ejecutado por Vagrant, actuando como un \textit{wrapper} simple para delegar la ejecución del proceso al orquestador principal con privilegios de \texttt{root}.
\subsection{Propósito}
El script es minimalista para simplificar la traza de logs de Vagrant, siendo su única tarea invocar al orquestador:
\begin{lstlisting}[language=bash, caption={Invocación del Orquestador}]
# Run root setup script
sudo bash -c 'sh /vagrant/2-Standby/scripts/primary/root\_setup.sh'
\end{lstlisting}
El uso de \texttt{sudo} es **necesario** ya que los pasos iniciales de instalación requieren privilegios de \texttt{root} para modificar el sistema operativo.
\section{Orquestador Principal (\texttt{root\_setup.sh})}

El script \texttt{root\_setup.sh} coordina la secuencia de 12 pasos (del 3 al 16) necesarios para construir un nodo primario de Oracle Data Guard completamente funcional, desde la preparación de discos hasta la instalación de \textit{cron jobs}.
\subsection{Secuencia de Ejecución}
La secuencia de ejecución NO es arbitraria; el orden es importante para las dependencias de Oracle:
\begin{enumerate}
    \item \textbf{Preparación de Discos}
    \item \textbf{Instalación de Paquetes OS}
    \item \textbf{Configuración de Red}
    \item \textbf{Configuración de Entorno Oracle}
    \item \textbf{Instalación de Software}
    \item \textbf{Creación de la Base de Datos con Data Guard}
    \item \textbf{Despliegue del Esquema de Aplicación}
    \item \textbf{Configuración de Automatización}
    \item \textbf{Configuración de Claves SSH}
    \item \textbf{Instalación de Cron Jobs}
\end{enumerate}
Si este script falla en cualquier punto, el despliegue se detiene, y se debe consultar el log completo de Vagrant (\texttt{03-primary-vagrant.log}) para identificar el paso fallido.

\section{Preparación de Discos (\texttt{prepare\_u01\_u02\_disks.sh})}

Este script configura el almacenamiento persistente para Oracle, siendo esencial para la estabilidad del sistema.
\subsection{Flujo de Configuración de Almacenamiento}
\begin{enumerate}
    \item \textbf{Particionado:} Utiliza \texttt{fdisk} en modo no interactivo para crear una partición primaria en \texttt{/dev/sdb} y \texttt{/dev/sdc}.
\item \textbf{Formateo:} Formatea las particiones como \textbf{XFS}, el sistema de archivos recomendado por Oracle para Linux por su rendimiento con \textit{datafiles} grandes.
\begin{lstlisting}[language=bash, caption={Comando de Formateo XFS}]
mkfs.xfs -f ${DISK_DEVICE}1
        \end{lstlisting}
    \item \textbf{Montaje Persistente:} Agrega entradas basadas en \textbf{UUID} a \texttt{/etc/fstab} para garantizar que los puntos de montaje (\texttt{/u01} y \texttt{/u02}) persistan después de reiniciar.
\end{enumerate}
\destacado{Resultado:} Se aíslan los datos y el software de Oracle (\texttt{/u01}) de los \textit{backups} y del \textit{Fast Recovery Area} (\texttt{/u02}).
\section{Instalación de Paquetes (\texttt{install\_os\_packages.sh})}

Este script prepara el sistema operativo a nivel de kernel y paquetes para la instalación de Oracle.
\subsection{Pasos de Configuración OS}
\begin{enumerate}
    \item \textbf{Pre-instalación Oracle:} Instala el paquete \texttt{oracle-database-preinstall-21c}, que automáticamente crea el usuario \texttt{oracle}, los grupos de privilegios (como \texttt{dba} y \texttt{dgdba}) y aplica límites de recursos.
\item \textbf{Paquetes Adicionales:} Instala utilidades como \texttt{unzip}, \texttt{rsync} (para transferencia de \textit{backups}), y \texttt{chrony} (para sincronización de tiempo).
\item \textbf{Parámetros del Kernel:} Configura límites de memoria compartida, semáforos, y \textbf{\textit{TCP tuning}} (memoria \texttt{rmem}/\texttt{wmem}) en \texttt{/etc/sysctl.d/98-oracle.conf}.
Este \textit{tuning} \textbf{afecta directamente el rendimiento} de la transferencia de \textit{archive logs} de Data Guard.
\item \textbf{Seguridad:} Deshabilita \texttt{SELinux} (a modo \texttt{permissive}) y \texttt{firewalld} para el entorno de desarrollo.
\end{enumerate}
El paquete \texttt{chrony} instalado es **un requisito** para la sincronización de tiempo, ya que Data Guard Broker la necesita para operar sin fallos de consistencia.

\section{Configuración de Red (\texttt{configure\_network.sh})}

La conectividad y la resolución de nombres son esenciales para la comunicación entre los nodos \textit{primary} y \textit{standby}.
\subsection{Configuración de Red Data Guard}
\begin{enumerate}
    \item \textbf{\texttt{/etc/hosts}:} Agrega las entradas de DNS local para \texttt{ol8-21-dg1.localdomain} y \texttt{ol8-21-dg2.localdomain} para asegurar que los nodos puedan encontrarse por nombre a través de la red \textit{Host-Only}.
\item \textbf{Hostname:} Configura el \textit{hostname} estático del nodo (ej., \texttt{ol8-21-dg1.localdomain}) usando \texttt{hostnamectl}. Oracle depende del \textit{hostname} correcto para registrar servicios.
\item \textbf{Sincronización de Tiempo (\texttt{chrony}):} Habilita y arranca el servicio \texttt{chronyd}.
La sincronización de tiempo evita errores de \textit{lag} y permite que Oracle confíe en los \textit{timestamps} de los \textit{redo logs}.
\begin{lstlisting}[language=bash, caption={Verificación de Sincronización}]
# Verificar que el servicio está activo
systemctl status chronyd

# Verificar la diferencia de tiempo
chronyc tracking
        \end{lstlisting}
\end{enumerate}
\advertencia{Si la configuración de \texttt{/etc/hosts} es incorrecta, la creación de la base de datos y la configuración del \textit{listener} fallarán, imposibilitando el inicio de Data Guard.}

\section{Entorno de Usuario Oracle (\texttt{oracle\_user\_environment\_setup.sh})}

Este script prepara el entorno de usuario \texttt{oracle} antes de la instalación del software.
\subsection{Configuración de Variables de Entorno}
\begin{enumerate}
    \item \textbf{Creación de \texttt{setEnv.sh}:} Genera el script que contiene todas las variables de Oracle (\texttt{ORACLE\_HOME}, \texttt{ORACLE\_SID}, \texttt{ORACLE\_UNQNAME}, \texttt{PATH}, \texttt{TNS\_ADMIN}).
\begin{lstlisting}[language=bash, caption={Variables Críticas de Entorno}]
export ORACLE_HOME=/u01/app/oracle/product/21.0.0/dbhome_1
export ORACLE_SID=cdb1
export ORACLE_UNQNAME=cdb1
export PATH=${ORACLE_HOME}/bin:${PATH}
export TNS_ADMIN=${ORACLE_HOME}/network/admin
        \end{lstlisting}
    \item \textbf{Persistencia:} Agrega el \textit{source} de \texttt{setEnv.sh} a \texttt{/home/oracle/.bash\_profile} para que las variables se carguen automáticamente en cada sesión.
\end{enumerate}
\destacado{Diferencia Clave:} En el \textit{primary}, \texttt{ORACLE\_UNQNAME} se establece como \texttt{cdb1}. En el \textit{standby}, el script equivalente lo establece como \texttt{cdb1\_stby}.
\section{Instalación de Software Oracle (\texttt{oracle\_db\_software\_installation.sh})}

Este es el paso de instalación de binarios, ejecutado como usuario \texttt{oracle} para asegurar permisos correctos.
\subsection{Instalación Silenciosa}
\begin{enumerate}
    \item \textbf{Copia y Descompresión:} El script primero copia el ZIP de Oracle de la carpeta compartida de Vagrant (\texttt{/vagrant\_software/}) a un directorio local (\texttt{/u01/software/}) y lo descomprime directamente en \texttt{ORACLE\_HOME}.
\item \textbf{\textit{Response File}:} Genera un \textit{response file} (\texttt{db\_install.rsp}) para la instalación silenciosa en modo \textbf{Software Only} (\texttt{INSTALL\_DB\_SWONLY}) y \textbf{Enterprise Edition} (\texttt{EE}).
\item \textbf{Ejecución del Instalador:} Ejecuta \texttt{runInstaller} en modo silencioso.
        \begin{lstlisting}[language=bash, caption={Ejecución runInstaller Silencioso}]
${ORACLE_HOME}/runInstaller -ignorePrereq -waitforcompletion -silent \
  -responseFile /tmp/db_install.rsp
        \end{lstlisting}
\end{enumerate}
Una vez que el \textit{runInstaller} finaliza, el orquestador (\texttt{root\_setup.sh}) ejecuta los scripts \texttt{orainstRoot.sh} y \texttt{root.sh} como \texttt{root} para completar la configuración de \textit{binarios} y \textit{permissions}.
\section{Creación de Base de Datos con Data Guard (\texttt{oracle\_create\_database.sh})}

Este es el \textbf{paso determinante} del despliegue del \textit{primary}, donde la base de datos se crea con la configuración necesaria para la replicación.
\subsection{Configuración de Red y DBCA}
\begin{enumerate}
    \item \textbf{Archivos TNS:} Crea \texttt{tnsnames.ora}, \texttt{listener.ora} y \texttt{sqlnet.ora}.
Este último configura \textbf{Network Encryption} (\texttt{AES256}) y \textit{checksums} (\texttt{SHA256}) para proteger los logs en tránsito.
\item \textbf{DBCA:} Crea la base de datos CDB (\texttt{cdb1}) con un PDB (\texttt{pdb1}) usando \texttt{dbca -silent} (8-12 minutos).
\item \textbf{Password File:} Crea el archivo \texttt{orapwcdb1} para autenticación remota del \textit{sysdba}, esencial para el \textit{Broker} de Data Guard.
\end{enumerate}

\subsection{Configuración Data Guard con DGMGRL}

El paso más importante es la preparación de la base de datos:
\begin{lstlisting}[language=bash, caption={Comando Moderno PREPARE DATABASE}]
dgmgrl / <<EOF
prepare database for data guard
  with db_unique_name is ${NODE1_DB_UNIQUE_NAME}
  db_recovery_file_dest is "${FRA_LOCATION}"
  db_recovery_file_dest_size is ${FRA_SIZE};
exit;
EOF
\end{lstlisting}
\destacado{Resultado:} Este comando habilita automáticamente el modo \texttt{ARCHIVELOG}, \texttt{FORCE LOGGING}, configura \texttt{STANDBY REDO LOGS}, \texttt{Flashback Database} y el \textit{Fast Recovery Area} (\texttt{FRA} en \texttt{/u02}), dejando el \textit{primary} 100\% listo para la conexión del \textit{standby}.
\section{Despliegue de Esquema (\texttt{deploy\_custom\_schema.sh})}

Una vez que la base de datos está preparada, este script despliega el esquema de aplicación.
\subsection{Flujo del Esquema}
\begin{enumerate}
    \item \textbf{Apertura de PDB:} Se asegura de que \texttt{pdb1} esté en modo \texttt{OPEN}.
\item \textbf{Creación de Usuario/Tablespace:} Ejecuta \texttt{Script\_SchemeTableSpace.sql} para crear el usuario (por defecto \texttt{HC}) y su \textit{tablespace}.
\item \textbf{Despliegue de Objetos:} Ejecuta \texttt{ScriptProyectoBD.SQL} para crear todos los objetos de la base de datos (tablas, secuencias, etc.). El contenido de este script es personalizable.
\item \textbf{Replicación:} Se fuerza un \textit{log switch} (\texttt{ALTER SYSTEM ARCHIVE LOG CURRENT}) para garantizar que todas estas sentencias DDL y DML de creación se empaqueten inmediatamente en un \textit{archive log} y estén listas para ser replicadas al \textit{standby}.
\end{enumerate}
El esquema de aplicación creado aquí es una de las primeras cosas que el \textit{standby} clonará. Si el \textit{schema} no se replica correctamente, Data Guard fallará en la validación.

\section{Configuración de Automatización (\texttt{setup\_automation.sh})}

Este script prepara los mecanismos de monitoreo y mantenimiento, desplegando los scripts de automatización.
\subsection{Despliegue de Scripts}
\begin{enumerate}
    \item \textbf{Copia de Scripts:} Los scripts de la carpeta compartida \texttt{/vagrant\_automation/} se copian a la ruta local del usuario \texttt{oracle}: \texttt{/home/oracle/automation/}.
\item \textbf{Permisos y Logs:} Se asignan permisos de ejecución (\texttt{chmod +x}) y se crean los directorios de \textit{logs} (\texttt{/home/oracle/logs/automation/}).
\end{enumerate}
\destacado{Scripts Desplegados (Monitoreo y Mantenimiento):}
\begin{itemize}
    \item \texttt{check\_log\_transport.sh} (para \textit{cron job} cada 10 min)
    \item \texttt{rman\_daily\_backup.sh} (para \textit{cron job} de backup diario)
    \texttt{delete\_archives\_primary.sh} (para \textit{cron job} de limpieza de logs)
    \item \texttt{verify\_dataguard.sh} (\textit{health check} manual completo)
\end{itemize}

\section{Configuración de Claves SSH para Backups (\texttt{setup\_ssh\_keys\_auto.sh})}

Este script habilita la transferencia segura y automatizada de \textit{backups} entre el \textit{primary} y el \textit{standby} sin intervención humana.
\subsection{Flujo de Configuración SSH}
\begin{enumerate}
    \item \textbf{Generación de Claves:} Genera un par de claves \texttt{RSA 4096-bit} para el usuario \texttt{oracle}.
Se genera sin \textit{passphrase} (\texttt{-N ""}) para permitir la automatización.
\begin{lstlisting}[language=bash, caption={Generación de Claves RSA}]
sudo -u oracle ssh-keygen -t rsa -b 4096 -f "${SSH_KEY}" -N "" \
  -C "oracle@${NODE1_FQ_HOSTNAME}"
        \end{lstlisting}
    \item \textbf{Exportación:} La clave pública (\texttt{id\_rsa.pub}) se copia a \texttt{/vagrant/ssh\_keys/primary\_id\_rsa.pub} para que el \textit{standby} pueda importarla en su archivo \texttt{authorized\_keys}.
\item \textbf{Configuración del Cliente:} Se configura \texttt{~/.ssh/config} con opciones para la conexión al \textit{standby}:
        \begin{itemize}
            \item \texttt{BatchMode yes} (no pide interacción).
\item \texttt{StrictHostKeyChecking accept-new} (no pide confirmación de \textit{host key}).
        \end{itemize}
\end{enumerate}
\destacado{Importancia para Data Guard:} Estas claves permiten que el \textit{cron job} \texttt{rman\_daily\_backup.sh} transfiera la copia del \textit{backup} diario al \textit{standby} a través de \texttt{rsync} de forma segura y desatendida.
\section{Instalación de Cron Jobs Automatizados (\texttt{install\_cron\_jobs.sh})}

El script detecta que es el nodo \textit{primary} y configura las siguientes tareas en el \textit{crontab} del usuario \texttt{oracle}:
\begin{enumerate}
    \item \textbf{Monitoreo de Replicación:} Cada 10 minutos (\texttt{*/10 * * * *}) se ejecuta \texttt{check\_log\_transport.sh}.
\item \textbf{Backup Diario:} Diariamente a la 1:00 AM (\texttt{0 1 * * *}) se ejecuta \texttt{rman\_daily\_backup.sh}.
\item \textbf{Limpieza de Archives:} Diariamente a las 2:00 AM (\texttt{0 2 * * *}) se ejecuta \texttt{delete\_archives\_primary.sh}.
\end{enumerate}
El orden de los \textit{cron jobs} es intencional: el \textit{backup} se realiza a la 1 AM, y la \textit{limpieza} se realiza una hora después (2 AM) para asegurar que la copia del \textit{backup} se haya completado antes de que se eliminen los \textit{archive logs} que este respalda.

\subsection{Despliegue Primario Completado}
El proceso de \textit{provisioning} del nodo primario está oficialmente completado.
\newpage

% *************************************************
% APÉNDICE B: FLUJO DE DESPLIEGUE DEL NODO STANDBY
% *************************************************

\chapter{Análisis Detallado: Despliegue del Nodo Standby}

\section{Punto de Entrada (\texttt{setup.sh})}

El script \texttt{setup.sh} es el primer script ejecutado por Vagrant en el nodo \textit{standby}.
\subsection{Propósito}
Actúa como un \textit{wrapper} simple para invocar el orquestador principal:
\begin{lstlisting}[language=bash, caption={Invocación del Orquestador Standby}]
# Run root setup script
sudo bash -c 'sh /vagrant/2-Standby/scripts/standby/root\_setup.sh'
\end{lstlisting}

\section{Orquestador Principal (\texttt{root\_setup.sh})}

El \texttt{root\_setup.sh} del \textit{standby} coordina 10 etapas para crear la réplica de Data Guard.
La duración total es menor que la del \textit{primary} ($\sim 20$-$30$ minutos).
\subsection{Secuencia de Ejecución}
La secuencia se adapta para la clonación de la base de datos:
\begin{enumerate}
    \item \textbf{Preparación de Discos}
    \item \textbf{Instalación de Paquetes OS}
    \item \textbf{Configuración de Red} (Hosts, Hostname, Chrony)
    \item \textbf{Configuración de Entorno Oracle} (con \texttt{ORACLE\_UNQNAME=cdb1\_stby})
    \item \textbf{Instalación de Software}
    \item \textbf{Creación del Standby} (\texttt{oracle\_create\_database.sh} con \textbf{RMAN DUPLICATE})
    \item \textbf{Configuración de Automatización}
    \item \textbf{Configuración de Claves SSH} (Importación de clave \textit{primary})
    \item \textbf{Instalación de Cron Jobs} (solo limpieza)
\end{enumerate}

\section{Creación del Standby Database (\texttt{oracle\_create\_database.sh})}

Esta etapa es el corazón del despliegue del \textit{standby} y la activación de la alta disponibilidad.
\subsection{RMAN DUPLICATE para Clonación}
\begin{enumerate}
    \item \textbf{Conectividad TNS:} Se configuran los archivos de red para permitir la conexión al \textit{primary} y a la propia instancia auxiliar.
\item \textbf{Instancia Auxiliar:} Se arranca una instancia mínima en modo \texttt{NOMOUNT}.
\item \textbf{EJECUCIÓN RMAN DUPLICATE:}
        \begin{lstlisting}[language=bash, caption={Comando Crítico de Clonación y Configuración}]
rman target sys/${SYS_PASSWORD}@cdb1.world auxiliary sys/${SYS_PASSWORD}@cdb1_stby.world <<EOF

duplicate target database
  for standby
  from active database
  dorecover
  spfile
    set db_unique_name='cdb1_stby'
  nofilenamecheck;
exit;
EOF
        \end{lstlisting}
\end{enumerate}
\destacado{Resultado:} Una copia idéntica de la base de datos \texttt{cdb1} (\texttt{CDB}, \texttt{PDB}, esquema de aplicación y datos) se crea en el almacenamiento local del \textit{standby}.
\subsection{Configuración Data Guard Broker y MAXAVAILABILITY}
\begin{enumerate}
    \item \textbf{Flashback Database:} Se habilita \texttt{FLASHBACK DATABASE} en el \textit{standby}.
\item \textbf{Data Guard Broker:} Se ejecuta \texttt{CREATE CONFIGURATION} y \texttt{ENABLE CONFIGURATION} para establecer la gestión.
\item \textbf{Modo de Protección:} Se establece \texttt{MAXAVAILABILITY} con \texttt{LogXptMode='FASTSYNC'} para garantizar \textbf{cero pérdida de datos}.
\item \textbf{Apply-ON:} Se inicia el proceso \texttt{MRP} (\textit{Managed Recovery Process}) para la aplicación continua de logs.
\end{enumerate}

\section{Configuración de Claves SSH y Cron Jobs}

\subsection{Flujo de Importación SSH}
\begin{enumerate}
    \item \textbf{Espera:} El script espera activamente la clave pública del \textit{primary} en \texttt{/vagrant/ssh\_keys/}.
\item \textbf{Importación:} Se copia la clave pública del \textit{primary} a \texttt{/home/oracle/.ssh/authorized\_keys} del \textit{standby}.
\item \textbf{Directorio de Backups:} Se crea \texttt{/u02/backup\_from\_primary} para recibir los \textit{backups} del \textit{primary} vía \texttt{rsync}.
\item \textbf{Exportación de Host Keys:} Se utiliza \texttt{ssh-keyscan} para exportar las \textit{host keys} del \textit{standby} al \textit{primary}.
\end{enumerate}
\destacado{Resultado:} El \textit{primary} puede ejecutar \texttt{rsync} al \textit{standby} sin autenticación interactiva.
\subsection{Tareas Programadas en Standby}
\begin{enumerate}
    \item \textbf{Limpieza de Archives Aplicados:} Diariamente a las 11:00 PM (\texttt{0 23 * * *}) se ejecuta \texttt{delete\_archives\_standby.sh}.
\end{enumerate}
El \textit{standby} solo tiene este \textit{cron job} de limpieza. El monitoreo, \textit{backup} y transferencia son responsabilidad exclusiva del \textit{primary}.

\newpage
% *************************************************
% APÉNDICE C: ANÁLISIS DE SCRIPTS DE AUTOMATIZACIÓN
% *************************************************

\chapter{Análisis Detallado: Scripts de Automatización}

Esta sección documenta los scripts de automatización ubicados en \texttt{/home/oracle/automation/} que se utilizan para mantener, monitorear y verificar la salud del sistema Data Guard de forma continua.
\section{Monitoreo de Transferencia (\texttt{check\_log\_transport.sh})}

\subsection{Propósito}
Verificar cada 10 minutos que los \textit{archive logs} se transfieren y aplican correctamente.
\subsection{Flujo del Script}
\begin{enumerate}
    \item \textbf{Consulta de Primary:} Obtiene \texttt{PRIMARY\_SEQ}.
    \item \textbf{Consulta de Standby:} Obtiene \texttt{STANDBY\_SEQ} (última aplicada).
\item \textbf{Cálculo de Gap:} \texttt{GAP = PRIMARY\_SEQ - STANDBY\_SEQ}.
    \item \textbf{Alerta:} Si \texttt{GAP} es mayor a $\mathbf{10}$ secuencias, registra una advertencia.
\item \textbf{Logging:} Registra el \textit{gap} y el \textit{apply lag} en \texttt{/home/oracle/logs/automation/log\_transport.log}.
\end{enumerate}

\section{Respaldo Diario RMAN (\texttt{rman\_daily\_backup.sh})}

\subsection{Proceso de Respaldo}
\begin{enumerate}
    \item \textbf{Configuración RMAN:} Establece \textit{retention policy} a \textbf{7 días} y \textit{deletion policy} a \texttt{APPLIED ON ALL STANDBY}.
\item \textbf{Backup Comprimido:} Ejecuta \texttt{BACKUP AS COMPRESSED BACKUPSET DATABASE PLUS ARCHIVELOG}.
\item \textbf{Transferencia:} Utiliza \texttt{rsync} sobre SSH para transferir la copia del \textit{backup} al \textit{standby} (\texttt{/u02/backup\_from\_primary/}).
\end{enumerate}

\section{Limpieza de Archives en Primary (\texttt{delete\_archives\_primary.sh})}

\subsection{Criterio de Eliminación}
Elimina \textit{archive logs} en el \textit{Primary} solo si han sido \texttt{APPLIED} en el \textit{Standby} y han sido \texttt{BACKED UP} (\textit{cron job} 2:00 AM).
\section{Limpieza de Archives en Standby (\texttt{delete\_archives\_standby.sh})}

\subsection{Criterio de Eliminación en Standby}
Elimina \textit{archive logs} en el \textit{Standby} solo si han sido \texttt{APPLIED} localmente \textbf{Y} tienen más de $\mathbf{3}$ días de antigüedad (\textit{cron job} 11:00 PM).
\section{Verificación de Configuración (\texttt{verify\_dataguard.sh})}

Herramienta de diagnóstico \textbf{manual} para el DBA.
\subsection{Reporte Integral}
Verifica \textbf{11} secciones que incluyen: \texttt{PRIMARY} status, \texttt{MRP} activo, \texttt{Apply Lag}, \texttt{Archive Gap} y \texttt{Broker Status} (\texttt{SUCCESS}).
\section{Scripts de Operación Manual}

\begin{itemize}
    \item \textbf{\texttt{manual\_sync.sh}:} Fuerza \texttt{ALTER SYSTEM ARCHIVE LOG CURRENT} para sincronización inmediata.
\item \textbf{\texttt{manual\_backup.sh}:} Ejecución de \textit{backup} RMAN a demanda.
    \item \textbf{\texttt{verify\_backup\_automation.sh}:} Confirma que \textbf{todos los componentes de automatización} (\textit{cron jobs}, \textit{scripts}, \textit{SSH keys}) están configurados y funcionales.
\end{itemize}

\newpage

% *************************************************
% APÉNDICE D: BIBLIOGRAFÍA
% *************************************************

\chapter{Bibliografía y Recursos Adicionales}
\addcontentsline{toc}{chapter}{Bibliografía y Recursos Adicionales}

\begin{thebibliography}{9}

\bibitem{oracle_dg_concepts}
Oracle Corporation. (2025). \textit{Oracle Data Guard Concepts and Administration 21c}.
\url{https://docs.oracle.com/en/database/oracle/oracle-database/21/sbydb/oracle-data-guard-concepts.html}

\bibitem{oracle_rman_dg}
Oracle Corporation. (2025). \textit{Using RMAN in Oracle Data Guard Configurations}.
\url{https://docs.oracle.com/en/database/oracle/oracle-database/21/sbydb/using-RMAN-in-oracle-data-guard-configurations.html}

\bibitem{oracle_rman_getting_started}
Oracle Corporation. (2025). \textit{Getting Started with RMAN}.
\url{https://docs.oracle.com/en/database/oracle/oracle-database/21/bradv/getting-started-rman.html}

\bibitem{oracle_xfs}
Oracle Corporation. (2024). \textit{XFS Data File System}.
\url{https://www.oracle.com/linux/technologies/xfs-overview.html}

\bibitem{oracle_base_broker}
Hall, T. (s.f.). \textit{Data Guard Setup Using Broker (21c)}. Oracle-Base.
\url{https://oracle-base.com/articles/21c/data-guard-setup-using-broker-21c}

\bibitem{oracle_base_prepare}
Hall, T. (s.f.). \textit{Data Guard: Prepare Database for Data Guard (21c)}. Oracle-Base.
\url{https://oracle-base.com/articles/21c/data-guard-prepare-database-for-data-guard-21c}

\bibitem{vagrant_provisioning}
HashiCorp. (2025). \textit{Vagrant Docs: Provisioning}.
\url{https://developer.hashicorp.com/vagrant/docs/provisioning}

\end{thebibliography}

\end{document}